%!TEX root = ../../springer-mono-format/book.tex
%!TEX root = ../../thesis-trier.tex

Here, we collect the order-theoretical notions that we need such 
as quasiorders, partial orders, the induced equivalence relation, upper sets, etc.
\Cref{prelim:gen:order:sec:def} give definitions and some basic equations. In \cref{prelim:gen:order:sec:eq-upper-lower},
we investigate upper and lower sets and the induced equivalence relation. Maximal and minimal elements are 
investigated in \cref{prelim:gen:order:sec:maximal-minimal},
and the set of maximal elements in \cref{prelim:gen:order:sec:set-maximal}.
\Cref{prelim:general:order:sec:seq-upper-sets} collects results on sequences of upper sets.
We later need all these notions when we introduce a special quasiorder on automata.

\subsection{Definitions and Basic Facts}
\label{prelim:gen:order:sec:def}

\begin{enumerate}
	\IN{quasiorder} 
	\index{preorder|see{quasiorder}}
	\item A \emph{quasiorder} (or \emph{preorder}) on a set $X$ is a reflexive and transitive relation. 
	We also write $(X, \lesssim)$ for $X$ together with the quasiorder $\lesssim$. 
	
	
	\item For a quasiordered set $(X, \lesssim)$
	we call two elements $x,y \in X$ \emph{incomparable} \IN{incomparable elements}
	if neither $x \lesssim y$ nor $y \lesssim x$ is true.
	
	\item Let $\lesssim$ be a quasiorder on $X$ and $S \subseteq X$.
	\begin{enumerate}   \IN{minimal element}
		\item We call $x \in X$ a \emph{minimal element} (resp.~\emph{maximal element}) 
		of~$S$ if $x \in S$ and for all $y \in S$, \IN{maximal element}
		if $y \lesssim x$ (resp.~$x \lesssim y$), then $x \lesssim y$ (resp.~$y \lesssim x$).
		
				\IN{maximum} \index{greatest element|see{maximum}}
				\item We call $x \in X$ a \emph{minimum}, or \emph{least element},
				(resp.~\emph{maximum}, or \emph{greatest element}) 
				of~$S$ if $x \in S$ 
				and for all $y \in S$ we have $x \lesssim y$ (resp.~$y \lesssim x$).
				\IN{minimum} \index{least element|see{minimum}}
		
		\IN{lower set!quasiorder} \IN{upper set!for quasiorder}
		\item The set $S$ is a \emph{lower set} (resp.~an \emph{upper set}) if $x \in S$ and 
		$y \lesssim x$ (resp.~$x \lesssim y$) imply $y \in S$ for $x,y \in X$. \IN{upper set!quasiorder} 
		
		\label{def:genupperset} \IN{generated!upper set} \IN{generated!lower set}
		\item The \emph{smallest lower set} (resp.~\emph{upper set}) \emph{containing $S$ (in $X$)}
		is $\lowerset{S} = \{ y \in X \mid \exists x \in S \colon y \lesssim x \}$
		(resp.~$\upperset{S} = \{ y \in X \mid \exists x \in S \colon x \lesssim y \}$).
	\end{enumerate}

\begin{center}
	\fbox{ \begin{minipage}{30em}			
			The notions of maximal/maximum and minimal/minimum are usually only defined for partially ordered
			sets, but here we use them for quasiordered sets as well.
		\end{minipage}
	}
\end{center}

	\IN{morphism!of quasiordered sets} 
	\index{monotone function|see{morphism}}
	
	\item %For sets $X, Y$ with quasiorders $\lesssim$ and $\precsim$ respectively,
	Given quasiordered sets $(X, \lesssim)$ and $(Y, \lessapprox)$
	an \emph{(order) morphism} (or \emph{monotone function}) is a function $\varphi \colon X \to Y$  
	such that, for all $x,y \in X$, we have
	\[
	x \lesssim y \Rightarrow \applyfunc{x}{\varphi} \lessapprox \applyfunc{y}{\varphi}.
	\] 
	Additionally, we call $\varphi \colon X \to Y$ an \emph{(order) isomorphism} if there exists another order morphism \IN{isomorphism!order isomorphism}
	$\psi \colon Y \to X$ such that $\concat{\psi}{\varphi} = \operatorname{id}_Y$
	and $\concat{\varphi}{\psi} = \operatorname{id}_X$.	
	
	\IN{partial order}
	\item A \emph{partial order} on a set $X$ is a quasiorder that is antisymmetric.
	
	\index{linear order|see{total order}}
	\IN{total order}
	\item A \emph{total} (or \emph{linear}) \emph{order} is a partial order $\le$ on $X$ such that,
	for all $x,y \in X$, we have
	$	x \le y \text{ or } y \le x.
	$
	
	\IN{topological sorting!of a partial order} \index{topological order|see{topological sorting}}
	\index{linear extension|see{topological sorting}}
	\item A \emph{topological sorting} (also \emph{topological order} or \emph{linear extension})
	of a partial order $\le$ on $X$ 
	is a total order $\le^*$ on $X$ such
	that the identity transformation $\operatorname{id} \colon X \to X$ is a morphism from $(X, \le)$ to $(X, \le^*)$.\footnote{According
		to Szpilrajn's extension theorem every partial order
		can be topologically sorted~\cite{szpilrajn1930extension}.}
		
	For finite~$X$, a topological sorting 
	can be identified
	with an indexing $x_1, \ldots, x_n$ %enumeration
	of all the elements from $X$ such that
	$x_i \le x_j$ implies~$i \le j$.
		


%    \IN{maximum} \index{greatest element|see{maximum}}
%    \IN{minimum} \index{least element|see{minimum}}
%    \item For a partially ordered set $(X, \le)$, we call $x \in X$ a \emph{minimum}, or \emph{least element},
%    (resp.~\emph{maximum}, or \emph{greatest element}) 
%    of~$S$ if $x \in S$ 
%    and for all $y \in S$ we have $x \le y$ (resp.~$y \le x$).
%    Due to antisymmetry, a minimum (resp.~maximum) is unique and, if it exists, the only minimal (resp.~maximal) element.
     

 \IN{equivalence relation induced!by quasiorder}
 \IN{induced equivalence relation|see {equivalence relation induced}}
 \item Given a quasiorder $\lesssim$ on a set $X$, the \emph{induced equivalence relation (by $\lesssim$)} is
 \begin{equation}\label{prelim:eqn:induced_eqrel}
 	x \simeq y \Leftrightarrow x \lesssim y \land y \lesssim x,
 \end{equation}
 and the well-defined \emph{induced partial order} $\le$
 on the equivalence classes $X / {\simeq}$ \IN{induced partial order}
 is defined by setting
 \begin{equation}\label{prelim:eqn:induced_pto}
 	[x]_{\simeq} \le [y]_{\simeq} \Leftrightarrow x \lesssim y.
 \end{equation}  
 By Equation~\eqref{prelim:eqn:induced_pto}
 the function $x \mapsto [x]_{\simeq}$ is a surjective order morphism.
 When the context is clear, we also denote
 the induced partial order on the equivalence classes by the same symbol~$\lesssim$.
\end{enumerate}

A minimum (resp.~maximum) is also a minimal (resp.~maximal) element. 
If not said otherwise, 
we mean minimal, maximal, minimum, and maximum w.r.t.~$S = X$. 
In a partially ordered set, a minimum (resp.~maximum) equals a unique minimal (resp.~maximal) element.
For $p \lesssim q$, we also say that $p$ is below~$q$, $q$ is above~$p$,
and that  these relations hold in a strict sense if additionally $q \not\lesssim p$.


\IN{component-wise order}
\begin{example}\label{prelim:general:order:ex:component-wise-order}
Let $(X, \lesssim)$ be a quasiordered set and $I$ be some set.
Then the set of functions~$X^I$ becomes a quasiordered set $(X^I, \lesssim^*)$
by setting $f \lesssim^*  g$ if and only if~$\applyfunc{x}{f} \lesssim \applyfunc{x}{g}$ for all $x \in I$
and~$f,g \in X^I$. If $\lesssim$ is a partial order, then $\lesssim^*$ is also a partial order,
called the \emph{component-wise order}.
Often, without special mentioning, 
we simply denote this quasiorder on~$X^I$ by~$\lesssim$ as well.
\end{example}
 
\begin{example}\label{prelim:ex:order:NzeroSigma-induced-pto}
	For some alphabet $\Sigma$, the sets of functions $\mathbb N_0^{\Sigma}$ and $[n]^{\Sigma}$
	are partially ordered
	by the component-wise order from \cref{prelim:general:order:ex:component-wise-order}
	derived from the usual order on $\mathbb N_0$.
\end{example}

\todo[inline]{erst die gleichungen, dann die aussagen?}

Let $(X, \lesssim)$ and $(Y, \lessapprox)$ be quasiordered sets 
with induced equivalence relations~$\simeq$ and~$\approxeq$, respectively,
and let $\varphi \colon X \to Y$ be an order morphism.

\begin{proposition} \label{prelim:prop:morphism_equiv_class}
	For $x \in X$ we have $\varphi([x]_{\simeq}) \subseteq [\varphi(x)]_{\approxeq}$,
	i.e., the morphism~$\varphi$ maps an equivalence class into an equivalence class.
	Furthermore, if $\varphi$ is an order isomorphism,
	we also have $\varphi([x]_{\simeq}) = [\varphi(x)]_{\approxeq}$.
\end{proposition}
\begin{proof}
	Let $y \in [x]_{\simeq}$, i.e., $x \simeq y$.
	Clearly, we have $\varphi(x) \approxeq \varphi(y)$.
	Thus $\varphi(y) \in [\varphi(x)]_{\approxeq}$.
	
	Next, suppose $\varphi$ is an order isomorphism.
	Then $\concat{\varphi}{\psi} = \operatorname{id}_X$
	for some order morphism $\psi \colon Y \to X$.
	Let $y \in [\varphi(x)]_{\approxeq}$.
	By \cref{prelim:prop:surj-inj-inverses}, 
	$\varphi$ is surjective.	
	Thus there exist $z \in X$ such that $y = \varphi(z)$.
	Since $\varphi(z) \approxeq \varphi(x)$,
	we find $z = \psi(\varphi(z)) \approxeq \psi(\varphi(x)) = x$.
	Thus $z \in [x]_{\simeq}$
	and so $y \in \varphi([x]_{\simeq})$.
\end{proof}

\begin{proposition} \label{prelim:prop:morphism_on_induced_pto} 
	%Let $(X, \lesssim)$ and $(Y, \lessapprox)$ be quasiordered sets and $\varphi \colon X \to Y$ be an order morphism.
	%If $\le$ and $\preceq$ respectively denote the induced partial orders,
	%If $\simeq$ and $\approxeq$ denote the induced equivalence relations on $X$ and $Y$, respectively,
	The morphism $\varphi$ induces a well-defined order morphism $\overline{\varphi} \colon X / {\simeq} \to Y / {\approxeq}$
	w.r.t.~the induced partial 
	orders given by $\applyfunc{[x]_{\simeq}}{\overline{\varphi}} = [\varphi(x)]_{\approxeq}$.
\end{proposition}
\begin{proof}
 Since $x \simeq y$ implies $\varphi(x) \simeq \varphi(y)$,
 $\overline{\varphi}$ is well-defined.
 By \cref{prelim:eqn:induced_pto} it is a morphism.	
\end{proof}

\begin{proposition} \label{prelim:general:order:prop:equiv-class-image-disjoint-or-subset}
	$\varphi(S) \cap T \ne \emptyset$ implies $\varphi(S) \subseteq T$ for equivalence classes $S \subseteq X$ and $T \subseteq Y$
	of the induced equivalence relation.
\end{proposition} 
\begin{proof}
		Suppose~$\varphi(S) \cap T \ne \emptyset$.
		Since all elements in~$\varphi(S)$ are equivalent w.r.t.~$\approxeq$,
		and~$T$, as an equivalence class, is a maximal set whose element are all equivalent,
		we find $\varphi(S) \subseteq T$.	
\end{proof}

The following example gives a  
bijective (order) morphism that is not an isomorphism.

\begin{example}\label{prelim:ex:order:bijmorph}
	Let $X = \{0,1\}$ and $\le$ be the partial order with $0 \le 1$. Then $\operatorname{id}_X$
	is a bijective morphism of $X$ ordered by equality to $X$ ordered by $\le$
	that is not an isomorphism.
\end{example}

In the following statements, $(X, \lesssim)$ is a quasiordered set with induced equivalence relation~$\simeq$
and induced partial order $\le$ on $X / {\simeq}$. 

For $x,y \in X$, the following hold:
\begin{align} 
	x \lesssim y & \Leftrightarrow y \in \,\upperset{\{x\}} 
	\Leftrightarrow \,\upperset{\{ y \}} \subseteq \,\upperset{\{ x \}}  \label{prelim:eqn:reach-upperset} \\
	x \simeq y & \Leftrightarrow \upperset{\{x\}} = \upperset{\{y\}} \label{prelim:general:order:eqn:equiv-uppersets} \\
\end{align}
For a morphism $\varphi \colon X \to Y$ %where $Y$ is quasiordered by $\lessapprox$
and $S \subseteq X$ we have:
\begin{align}
	\varphi(\upperset{S}) & \subseteq \upperset{\varphi(S)} \label{prelim:general:order:eqn:morphism-upperset} \\
	\text{$S$ is an upper set} & \Rightarrow \text{$\varphi^{-1}(S)$ is an upper set}               \label{prelim:general:order:eqn:morphism-upperset-inverse}
\end{align} 
For $S, T \subseteq X$ we have:
\begin{align}
	\text{$S, T$ are upper sets} & \Rightarrow \text{$S \cup T$ and $S \cap T$ are upper sets} \label{prelim:general:order:eqn:upper-union-int}\\
	\text{$S, T$ are disjoint upper sets} & \Rightarrow \text{$\forall x \in S \ \forall y \in T \colon x \not\lesssim y \land y \not\lesssim x$} \label{prelim:general:order:eqn:disjoint-upper-sets} \\
	\text{$S$ is an upper set} & \Leftrightarrow S = \bigcup_{x \in S} \upperset{\{x\}} \label{prelim:general:order:eqn:upper-set} \\
	\text{$S$ is an upper set} & \Leftrightarrow \text{$X \setminus S$ is a lower set} \label{prelim:order:prop:upper-complement-lower} 
\end{align}
Similar statements to 
\cref{prelim:eqn:reach-upperset,prelim:general:order:eqn:equiv-uppersets,prelim:general:order:eqn:morphism-upperset,prelim:general:order:eqn:upper-union-int,prelim:general:order:eqn:disjoint-upper-sets,prelim:general:order:eqn:upper-set} 
also hold for lower sets. 

\begin{remark}\label{prelim:general:order:rem:morphism-upperset}
	The inclusion from \cref{prelim:general:order:eqn:morphism-upperset} 
	can be proper. For example, let $X = \{x_1, x_2\}$ where $x_1 \lesssim x_2$.
	Then with the morphism $\varphi \colon X \to X$ defined by $\varphi(x_1) = \varphi(x_2) = x_1$
	we have $\varphi(\upperset{\{x_1\}}) = \{x_1\}$
	and $\upperset{\{\varphi(x_1)\}} = \{x_1, x_2\}$. This also shows
	that a morphism does not necessarily map upper sets to upper sets.
	More generally, if all elements are comparable in an arbitrary set~$X$ with at least two elements,
	then any constant map is an order morphism
	such that the inclusion from \cref{prelim:general:order:eqn:morphism-upperset}
	is proper.
\end{remark}

\begin{remark}
	Let $x \in X$. Then $\upperset{\{x\}} \setminus [x]_{\simeq}$ is a upper set
	and $\lowerset{\{x\}} \setminus [x]_{\simeq}$ is a lower set.
\end{remark}

From \cref{prelim:order:prop:upper-complement-lower} 
we obtain the next statement.

\begin{proposition} \label{prelim:general:order:max-proper-upper-iff-min-nonempty-lower}
	$S \subseteq X$ is a maximal proper upper set w.r.t.~set inclusion
	if and only if $X \setminus S$ is a minimal non-empty lower set w.r.t.~set inclusion.
\end{proposition}

\begin{remark}
	\label{prelim:general:order:rem:upper-and-lower}
	In most of the remaining statements
	of the present section that make a claim on upper and lower sets, 
	we only give the proofs for upper sets, since the corresponding statement for lower sets
	follows similarly.
\end{remark}

\subsection{Equivalence Classes and Upper and Lower Sets}
\label{prelim:gen:order:sec:eq-upper-lower}

Here, \cref{prelim:general:order:subset-in-equiv}
relates upper and lower sets to equivalence classes,
and we list consequences in \cref{prelim:general:order:upper-union-equiv-class,prelim:automata:cor:SsetminusT}.
\Cref{prelim:automata:prop:TsetminusS-strong-conn,prelim:general:order:lem:diff-in-equiv-class,prelim:general:order:lem:upper-lower-diff-maximal} give relations between maximal and minimal upper and lower sets
and equivalence classes.

In this section, let $(X, \lesssim)$ be a quasiorder and let~$\simeq$
be its induced equivalence relation.

\begin{proposition} \label{prelim:general:order:subset-in-equiv}
	For $S \subseteq X$, the following are equivalent.
	\begin{enumerate}
		\item \label{prelim:general:order:subset-in-equiv:1} $S$ is contained in an equivalence class for~$\simeq$.
		
		\item \label{prelim:general:order:subset-in-equiv:2} For every upper set~$T \subseteq X$
		we have $S \subseteq T$
		or $S \cap T = \emptyset$. 
		
		\item \label{prelim:general:order:subset-in-equiv:3} For every lower set~$T \subseteq X$
		we have $S \subseteq T$
		or $S \cap T = \emptyset$. 		
	\end{enumerate}
\end{proposition}
\begin{proof}
$\ref{prelim:general:order:subset-in-equiv:1} \Rightarrow~\ref{prelim:general:order:subset-in-equiv:2}$: Let $T \subseteq X$ be an upper set.
 Suppose we have $x \in S \cap T$.
 Since~$T$ is an upper set, we have $\upperset{\{x\}} \subseteq T$.
 By \cref{prelim:eqn:reach-upperset},
 we find $[x]_{\simeq} \subseteq \upperset{\{x\}}$.
 Thus $S\subseteq T$ since $S \subseteq [x]_{\simeq}$ by assumption.
 
 	
$\ref{prelim:general:order:subset-in-equiv:2} \Rightarrow~\ref{prelim:general:order:subset-in-equiv:3}$: Let $T \subseteq X$ be a lower set. Since $X \setminus T$ is an upper set by \cref{prelim:order:prop:upper-complement-lower}, by \cref{prelim:general:order:subset-in-equiv:2}, $S$ is either contained in $X \setminus T$ or $T$.
 
$\ref{prelim:general:order:subset-in-equiv:3} \Rightarrow~\ref{prelim:general:order:subset-in-equiv:1}$: Let $x, y \in S$. 
 By assumption, we find $S \subseteq \,\lowerset{\{y\}}$
 and $S \subseteq \,\lowerset{\{x\}}$.
 This gives $x \lesssim y$ and $y \lesssim x$, respectively.
 Thus $x \simeq y$ and so $S$ is contained in an equivalence class. 
% Suppose $S$ is not contained in an equivalence class. 
% Then there exist $p, q \in S$
% such that $p \notin \upperset{\{q\}}$.
% Hence, $p \in S \setminus T$
% and $q \in S \cap T$ for the upper set 
% $T = \upperset{\{q\}}$.
\end{proof}

\begin{corollary} \label{prelim:general:order:upper-union-equiv-class}
	Every upper or lower set $S \subseteq X$ is a union of equivalence classes.
\end{corollary}
\begin{proof}
 Clearly, $S \subseteq \bigcup_{y \in S} [y]_{\simeq}$.
 By \cref{prelim:general:order:subset-in-equiv}, 
 we have $\bigcup_{y \in S} [y]_{\simeq} \subseteq S$.
\end{proof}

\begin{corollary} \label{prelim:automata:cor:SsetminusT}
	Let $S, T \subseteq X$ be both upper or lower sets.	
	If $T \setminus S$ is non-empty and contained in an equivalence class,
	then it equals this equivalence class.
\end{corollary}
\begin{proof}
    Without loss of generality, suppose $S$ and~$T$ are upper sets.
	Let $x \in T \setminus S$.
	By assumption, $T \setminus S \subseteq [x]_{\simeq}$.
	By \cref{prelim:general:order:subset-in-equiv},
	$[x]_{\simeq} \subseteq T$
	and $[x]_{\simeq} \cap S = \emptyset$.
	Therefore, $[x]_{\simeq} \subseteq T \setminus S$.
\end{proof}

\begin{lemma}\label{prelim:general:order:lem:diff-in-equiv-class}
	Let $S,T \subseteq X$.
	If $T \setminus S$ is contained in an equivalence class, then
	for all upper sets $R \subseteq X$ with $S \cap T \subseteq R \subseteq T$, 
	we have $R = S \cap T$ or $R = T$.
\end{lemma}
\begin{proof} 
	Let~$R$ be an upper set	with $S\cap T \subseteq R \subseteq T$.
	By \cref{prelim:general:order:subset-in-equiv}
	either we have $( T \setminus S ) \subseteq R$, which gives $T = (S \cap T) \cup (T \setminus S) \subseteq R$
	and so $R = T$,
	or $(T \setminus S) \cap R = \emptyset$,
	which gives 
	$
	R = R \cap T = R \cap ( (S \cap T) \cup (T \setminus S) ) = ( R \cap ( S \cap T ) ) \cup ( R \cap ( T \setminus S ) )
	= R \cap (S \cap T) = S \cap T.
	$		
\end{proof}

\begin{lemma}\label{prelim:general:order:lem:upper-lower-diff-maximal}
	Let $S,T \subseteq X$.
	Then the following are equivalent.
\begin{enumerate}
	  \item \label{prelim:general:order:lem:upper-lower-diff-maximal:1} For all upper sets $R \subseteq X$ with $S \cap T \subseteq R \subseteq T$, we have $R = S \cap T$ or $R = T$.
	
	\item \label{prelim:general:order:lem:upper-lower-diff-maximal:2} For all lower sets $R \subseteq X$ with $S \cap T \subseteq R \subseteq T$, 
	we have $R = S \cap T$ or $R = T$.    	
\end{enumerate}	 	
\end{lemma}
\begin{proof}
	Let~$R$ be a lower (resp.~an upper) set
	with $S\cap T \subseteq R \subseteq T$.
	Then $X \setminus T \subseteq X \setminus R \subseteq X \setminus (S \cap T) = (X \setminus S) \cup (X \setminus T)$.
	Since $X \setminus R, X\setminus S$, $X \setminus T$,
	and $( X\setminus T ) \cup (X \setminus T)$ 
	are upper (resp.~lower) sets by \cref{prelim:order:prop:upper-complement-lower,prelim:general:order:eqn:upper-union-int},
	from \cref{prelim:general:order:lem:upper-lower-diff-maximal:1}
	we obtain $X \setminus R = X \setminus T$
	or $X \setminus R = X \setminus ( S \cap T )$,
	and so $R = T$ or $R = S \cap T$.
    This proves \cref{prelim:general:order:lem:upper-lower-diff-maximal:2}.
    The converse implication follows similarly.
 \end{proof}

An analogous equivalence to \cref{prelim:general:order:lem:upper-lower-diff-maximal}, between upper and lower sets but with $S \subseteq R \subseteq S \cup T$ in place of $S \cap T \subseteq R \subseteq T$, can be established. We omit such a statement here.

\begin{proposition} \label{prelim:automata:prop:TsetminusS-strong-conn}
	Let $S, T \subseteq X$ be both upper sets or both lower sets.
    Then the following are equivalent.
    
    \begin{enumerate}
    \item \label{prelim:automata:prop:TsetminusS-strong-conn:1} $T \setminus S$
    is contained in an equivalence class.
    
    \item \label{prelim:automata:prop:TsetminusS-strong-conn:2} For all upper sets $R \subseteq X$ with 
    $S \cap T \subseteq R \subseteq T$ we have $R = S \cap T$ or $R = T$.
    \end{enumerate}	
\end{proposition}
\begin{proof}
$\ref{prelim:automata:prop:TsetminusS-strong-conn:1} \Rightarrow \ref{prelim:automata:prop:TsetminusS-strong-conn:2}$: This follows from \cref{prelim:general:order:lem:diff-in-equiv-class} (note that this statement applies to any sets, not just upper/lower sets).
		
$\ref{prelim:automata:prop:TsetminusS-strong-conn:2} \Rightarrow \ref{prelim:automata:prop:TsetminusS-strong-conn:1}$: 
First, suppose $S$ and $T$ are upper sets.
Suppose $T \setminus S$ is not contained in an equivalence class.
Then we have $x \not\lesssim y$ for some $x,y \in T \setminus S$.
Let $R = \upperset{\{x\}}$. We have $y \notin R$.
By \cref{prelim:general:order:eqn:upper-union-int}, $(S \cap T) \cup (R \cap T)$ is an upper set.
This upper set properly contains~$S\cap T$
and is properly contained in~$T$.
%
% Alternativer Beweis
%First, suppose $S$ and $T$ are upper sets. Let $R \subseteq X$
%be an upper set.
%Set $K = ( S \cap T ) \cup ( R \cap T )$.
%By \cref{prelim:general:order:eqn:upper-union-int}, $K$ is an upper set.
%Since $S\cap T \subseteq K \subseteq T$, from the assumption
%we get $K = S \cap T$ or $K = T$.
%In the former case we have $(T \setminus S) \cap R = \emptyset$,
%and in the latter case we have $(T \setminus S) \subseteq R$.
%By \cref{prelim:general:order:subset-in-equiv}, we find that $T \setminus S$
%is contained in an equivalence class.

Next, if~$S$ and $T$ are lower sets, we use the same reasoning
but with \crefitem{prelim:general:order:lem:upper-lower-diff-maximal}{prelim:general:order:lem:upper-lower-diff-maximal:2}
and lower set instead of upper sets.\footnote{Alternatively, for a lower set $R \subseteq X$
consider the lower set $K = (S \cup T) \cap (R \cup S)$.
Since $S \subseteq K \subseteq S \cup T$,
we have $( X \setminus S ) \cap ( X \setminus T ) \subseteq X \setminus K \subseteq X \setminus S$
and as these sets are upper sets by \cref{prelim:general:order:eqn:upper-union-int},
from the assumption we get $X \setminus K = X \setminus ( S \cup T )$
or $X \setminus K = X \setminus S$, equivalently $K = S \cup T$ or $K = S$. In the former case, $T \setminus S \subseteq R$,
and in the latter case, $(T \setminus S) \cap R = \emptyset$.
So, from \cref{prelim:general:order:subset-in-equiv}, we get that~$T \setminus S$ is contained in an equivalence class.}
\end{proof}

By \cref{prelim:automata:cor:SsetminusT,prelim:automata:prop:TsetminusS-strong-conn}
for upper sets $S,T$, if $T \setminus S$ is non-empty and contained in an equivalence class,
then $T = (S \cap T) \cup [x]_{\simeq}$ for some $x \in T \setminus S$.


\subsection{Maximal and Minimal Elements}
\label{prelim:gen:order:sec:maximal-minimal}

\Cref{prelim:general:order:cor:min-equiv-class}
characterizes upper and lower sets that are equivalence classes.
In \cref{prelim:general:order:XS-only-uppersets}
we characterize uniqueness of a proper non-empty upper or lower set.
\Cref{prelim:order:prop:num-disj-upper,prelim:order:cor:num-disj-upper}
gives quantitative statements on the number of disjoint upper and lower sets,
as well as on incomparable maximal and minimal elements in finite sets.
Lastly, in \cref{prelim:order:prop:maximum,prelim:order:prop:maximum2,prelim:general:order:prop:unique-equiv-and-upper}
we list results concerning on maxima and minima.

In this section, let $(X, \lesssim)$ be a quasiorder and let~$\simeq$ be its induced equivalence relation.

\begin{proposition} \label{prelim:general:order:cor:min-equiv-class} \label{prelim:general:order:prop:min-equiv-class} \label{prelim:prop:sink-comp}
	The following are equivalent for $S \subseteq X$:
	
	\begin{enumerate}
	\item \label{prelim:general:order:cor:min-equiv-class:1} $S$ is a minimal non-empty upper (resp.~lower) set w.r.t.~set inclusion.
	
	\item \label{prelim:general:order:cor:min-equiv-class:2} $S$ is an equivalence class that is an upper (resp.~a lower) set.
	
	\item \label{prelim:general:order:cor:min-equiv-class:2a} $S$ is non-empty such that $S = \bigcap_{x \in S} \upperset{\{x\}}$ (resp.~$S = \bigcap_{x\in S} \lowerset{\{x\}}$).
		
	\item \label{prelim:general:order:cor:min-equiv-class:3} \label{prelim:eqn:sink-comp} $S$ is non-empty such that $S = \upperset{\{x\}}$ (resp.~$S = \,\lowerset{\{x\}}$) for each $x \in S$.
		
	\item \label{prelim:general:order:cor:min-equiv-class:4} 
	Each~$x \in S$ is maximal (resp.~minimal) and $S = [x]_{\simeq}$.
	
	\item \label{prelim:general:order:cor:min-equiv-class:4b} We have $S = [x]_{\simeq}$ for some maximal (resp.~minimal) $x \in X$.
				
	\item \label{prelim:general:order:cor:min-equiv-class:5} $S$ is a maximal (resp.~minimal) equivalence class in $(X / {\simeq}, \le)$.	
	\end{enumerate}
\end{proposition}
\begin{proof}
	As mentioned in \cref{prelim:general:order:rem:upper-and-lower}, we prove the result only for upper sets.
	
$\ref{prelim:general:order:cor:min-equiv-class:1}\Rightarrow~\ref{prelim:general:order:cor:min-equiv-class:2}$: By minimality, there exists no upper set~$R \subseteq X$ with $R \subseteq S$
 such that $R \ne \emptyset$ and $R \ne S$.
 Hence, by \cref{prelim:automata:prop:TsetminusS-strong-conn} applied with $\emptyset$ and~$S$, 
 $S$ is contained in an equivalence class.
 Since~$S \ne \emptyset$, by \cref{prelim:automata:cor:SsetminusT} (or by \cref{prelim:general:order:upper-union-equiv-class}), 
 it equals an equivalence class.

$\ref{prelim:general:order:cor:min-equiv-class:2}\Rightarrow~\ref{prelim:general:order:cor:min-equiv-class:2a}$
 Let $x \in S$. Since~$S$ is an equivalence class we have $S \subseteq \upperset{\{x\}}$
 by \cref{prelim:general:order:upper-union-equiv-class}.
 Thus $S \subseteq \bigcap_{x \in S} \upperset{\{x\}}$.
 Since~$S$ is an upper set, we have $\upperset{\{x\}} \subseteq S$,
 Thus $\bigcap_{x \in S} \upperset{\{x\}} \subseteq S$. Since~$S$ 
 is an equivalence class it is non-empty.
 
$\ref{prelim:general:order:cor:min-equiv-class:2a}\Rightarrow~\ref{prelim:general:order:cor:min-equiv-class:3}$
 Let $x \in S$. Then $S = \bigcap_{y \in S}\upperset{\{y\}} \subseteq \upperset{\{x\}}$.
 It is easy to see that $S = \bigcap_{y \in S} \upperset{\{y\}}$ is an upper
 set. Thus we have $\upperset{\{x\}} \subseteq S$ for $x \in S$
 since $\upperset{\{x\}}$ is the smallest upper set generated by~$x$.  

$\ref{prelim:general:order:cor:min-equiv-class:3}\Rightarrow~\ref{prelim:general:order:cor:min-equiv-class:4}$: Let $x \in S$ and assume $x \lesssim y$ for some $y \in X$. Then $y \in \upperset{\{x\}} = S$.
By assumption, $\upperset{\{y\}} = S = \upperset{\{x\}}$.
Thus, with \cref{prelim:general:order:eqn:equiv-uppersets}, we get $x \simeq y$.
So~$x$ is maximal.
Next, let $z \in [x]_{\simeq}$.
Then $z \in \upperset{\{x\}} = S$. Thus we have $[x]_{\simeq} \subseteq S$.
Now, suppose $z \in S$. Then $\upperset{\{z\}} = S = \upperset{\{x\}}$
and so $z \simeq x$ by \cref{prelim:general:order:eqn:equiv-uppersets}.
This gives $[x]_{\simeq} \subseteq S$.

$\ref{prelim:general:order:cor:min-equiv-class:4}\Rightarrow~\ref{prelim:general:order:cor:min-equiv-class:4b}$: This is clear.

$\ref{prelim:general:order:cor:min-equiv-class:4b}\Rightarrow~\ref{prelim:general:order:cor:min-equiv-class:5}$:
 Let $S = [x]_{\simeq}$ with $x$ maximal in~$X$.
 Suppose $[x]_{\simeq} \le [y]_{\simeq}$ for some $y \in X$.
 Then we have $x \lesssim y$. Since~$x$ is maximal,
 we find $x \simeq y$ and so $[x]_{\simeq} = [y]_{\simeq}$. Thus $[x]_{\simeq}$
 is maximal in $(X / {\simeq}, \le)$.

$\ref{prelim:general:order:cor:min-equiv-class:5}\Rightarrow~\ref{prelim:general:order:cor:min-equiv-class:1}$: 
 Let $x \in S$ and $y \in X$ with $x \lesssim y$.
 Since~$S$ is maximal, we find $x \simeq y$ and so $x \in S$. Hence~$S$
 is an upper set. 
 Next, let~$R \subseteq X$ be a non-empty upper set
 with $R \subseteq S$. Since~$S$ is an equivalence class,
 \cref{prelim:general:order:upper-union-equiv-class} (or \cref{prelim:general:order:lem:diff-in-equiv-class})
 implies that $S \subseteq R$. Therefore,~$S$ is minimal w.r.t.~set inclusion
 among all non-empty upper sets.
\end{proof}

\begin{corollary}
	An element $x \in X$ is maximal if and only if $[x]_{\simeq}$ is an upper set.	
\end{corollary}

\begin{corollary}
	If~$X$ is finite, then every upper set contains
	an upper upper set forming a maximal equivalence class in $(X / {\simeq}, \le)$.
\end{corollary}

\begin{corollary}
	An upper set is a maximal proper set w.r.t.~set inclusion
	if and only if its complement is a minimal equivalence class
	for the induced partial order.
\end{corollary}
\begin{proof} This follows from directly from \cref{prelim:general:order:max-proper-upper-iff-min-nonempty-lower,prelim:general:order:cor:min-equiv-class}.
\end{proof}

\begin{proposition} \label{prelim:general:order:XS-only-uppersets}
	Let $S \subseteq X$ be a non-empty and proper upper (resp.~lower) set.
	Then $S$ and~$X$ are the only non-empty upper (resp.~lower) sets if and only
	if~$S$ and $X \setminus S$ are equivalence classes and~$X \setminus S$ is not an upper (resp.~a lower) set.
\end{proposition}
\begin{proof}
	Suppose~$S$ and~$X$ are the only non-empty upper sets.
	By considering $T = X$ and~$S$
	in  \cref{prelim:automata:prop:TsetminusS-strong-conn},
	we find that $X \setminus S$ is contained in an equivalence class.
	From \cref{prelim:automata:cor:SsetminusT} 
	we get that it equals that equivalence class.
	Next, for $x \in S$, since  $\upperset{\{x\}} \subseteq S$,
	we get $\upperset{\{x\}} = S$
	and so~$S$ is an equivalence class by \cref{prelim:general:order:cor:min-equiv-class}.
	By assumption, $X \setminus S$ is not an upper set.	
		
	Conversely, suppose~$S$ and $X \setminus S$ are equivalence
	classes. Since all upper sets are unions of equivalence classes and~$X$ is non-empty,
	$S$, $X \setminus S$, and~$X$ are all possible non-empty upper sets.
	So, if $X \setminus S$ is not an upper set, only $S$ and~$X$ remain as upper sets.
\end{proof}

\begin{proposition} \label{prelim:order:prop:num-disj-upper}
	Suppose~$X$ is finite. Then the following numbers are equal:
	
	\begin{enumerate}
	\item The maximal number of pairwise disjoint non-empty upper (resp.~lower) sets.
	
	\item The number of maximal (resp.~minimal) elements in $(X / {\simeq}, \le)$.
	
	\item The maximal number of incomparable maximal (resp.~minimal) element in~$X$.
	\end{enumerate}	
\end{proposition}
\begin{proof}
	Let $m_1, m_2, m_3$ be the maximal number of pairwise disjoint non-empty upper sets,
	of distinct maximal elements in $(X / {\simeq}, \le)$, 
	and the maximal number of incomparable maximal elements in~$X$, respectively.	

By \cref{prelim:general:order:cor:min-equiv-class} maximal elements in $X / {\simeq}$
are upper sets that are equivalence classes, hence pairwise disjoint.
This gives $m_1 \ge m_2$.	

By the definition of the induced partial order,
incomparable maximal elements induce
distinct maximal elements in $X / {\simeq}$, which gives $m_2 \ge m_3$.

By \cref{prelim:general:order:eqn:disjoint-upper-sets}, we have $m_3 \ge m_1$.	
\end{proof}

\begin{corollary}  \label{prelim:order:cor:num-disj-upper}
	For finite~$X$, there are at most~$k$ distinct maximal (resp.~minimal) elements in the partially ordered set 
	$(X /{\simeq}, \le)$ if and only if, among $k+1$ upper (resp.~lower) sets, two have a non-empty intersection. 
\end{corollary}

The following statements are easy to verify and we omit their proofs.

\begin{proposition} \label{prelim:order:prop:maximum}
	The following are equivalent for $x \in X$:	
	\begin{enumerate}
		\item $x \in X$ is a maximum (resp.~minimum) for $X$.
		\item $[x]_{\simeq}$ is a maximum (resp.~minimum) in $(X / {\simeq}, \le)$. 
		\item $[x]_{\simeq} \subseteq S$ for every upper (resp.~lower) set $S \subseteq X$. % da vereingung, und wenn das nicht gelten würden, kann man größtes [x] [x] darin finden, und die gehören zu verschiedenen maximalen
	\end{enumerate} 	
\end{proposition}

\begin{proposition} \label{prelim:order:prop:maximum2}
	The following are equivalent:
	\begin{enumerate}		
		\item There exists a maximum  (resp.~minimum) in $X$ for $\lesssim$.
		
		\item For $x,y \in X$ there exists $z \in X$ with $x \lesssim z$ and~$y \lesssim z$ (resp.~$z \lesssim x$ and $z \lesssim y$). 
		
		\item For two maximal (resp.~minimal) $x,y \in X$ we have $x \simeq y$.
		
		\item For every two upper (resp.~lower) sets $S, T \subseteq X$ we have $S \cap T \ne \emptyset$.		
		
	\end{enumerate} 	
\end{proposition}

\begin{proposition}\label{prelim:general:order:prop:unique-equiv-and-upper}
	Let $S \subseteq X$ be an equivalence class.
	If~$S$ is a maximum (resp.~minimum) in~$X / {\simeq}$,
	then it is the only upper (resp.~lower) set among the equivalence classes.
	Furthermore, if~$X$ is finite and~$S$ the only upper (resp.~lower) set that is an equivalence
	class, then it is a maximum (resp.~minimum) in $X / {\simeq}$.
\end{proposition}
\begin{proof}
	We only give the proof for upper sets and maxima.
	First, suppose~$S$ is an equivalence class that is a maximum in $X / {\simeq}$.
	Let $T \subseteq X$ be an equivalence class that is an upper set.
	By \cref{prelim:order:prop:maximum}, we have $S \subseteq T$.
	Since~$S$ and~$T$ are equivalence classes,
	we find $S = T$.
	By \cref{prelim:general:order:cor:min-equiv-class}, $S$ is
	an upper set, and so a unique equivalence class
	that is an upper set.
	
	Conversely, suppose~$X$ is finite and $S$ is 
	a unique upper set that is an equivalence class.
	By \cref{prelim:general:order:cor:min-equiv-class}, $S$
	is maximal in $X / {\simeq}$ and unique with this property.
	By finiteness, for each $x \in X$
	there exists a maximal $y \in X$ such that $x \le y$.
	Since~$S$ is the unique maximal element,
	we find $y \in S$.
	This yields that $S$ is a maximum in $X / {\simeq}$.	
\end{proof}

To show that the converse implication
from \cref{prelim:general:order:prop:unique-equiv-and-upper}
might fail for infinite~$X$, consider $X = \mathbb N_0 \cup \{ a \}$
with the usual order on~$\mathbb N_0$
and~$a$ being incomparable to any other element.
Then $\{a\}$ is a unique upper set that is an equivalence class, but it is not a maximum
w.r.t.~the induced partial order on the equivalence classes (which results in the same order structure
since every equivalence class is a singleton).


\subsection{The Set of Maximal Elements}
\label{prelim:gen:order:sec:set-maximal}

In this section, let $(X, \lesssim)$ and
\[
S = \{ x \in X \mid \text{$x$ is maximal}\}
\]
be the set of maximal elements. We denote the induced equivalence with~$\simeq$.

\begin{proposition} \label{prelim:order:maximal-elements-upper-set}
	The set of maximal elements~$S$ is an upper set and we have
	\begin{equation}\label{prelim:general:order:set-of-maximal-elements}
		S = \bigcup_{\text{$x$ is maximal}} [x]_{\simeq}.
	\end{equation}
\end{proposition}
\begin{proof}
	\Cref{prelim:general:order:set-of-maximal-elements} is obvious,
	and together with \cref{prelim:general:order:cor:min-equiv-class}
	we obtain that~$S$ is an upper set.	
\end{proof}

Next, we investigate~$S$ for finite~$X$. Similar statements could be obtained
for the set of minimal elements. But since we only need the set of maximal elements
later, we do not formulate corresponding statements for this set.

\begin{lemma} \label{prelim:general:order:lem:finite-maximal}
	If~$X$ is finite, then
	for each $x \in X$ there exists $y \in S$ such that $x \lesssim y$.
\end{lemma}
\begin{proof}
	Assume there exists $x \in X$ for which this condition is false.
	Set $x_0 = x$ and, inductively, for $x_{i}$ with $i \ge 0$
	there exists $x_{i+1} \in X$ such that $x_i \lesssim x_{i+1}$
	but $x_{i+1} \not\lesssim x_i$,
	since otherwise $x_i$ would be a maximum
	in~$\{x_0, x_1, \ldots, x_i\}$
	and so also maximal such that $x \lesssim x_i$.
	Now, in this way, we obtain an infinite and strictly increasing sequence,
	which is not possible for finite~$X$.
\end{proof}

\begin{proposition} \label{prelim:prop:maximal-elements-and-upper-sets}
	Suppose~$X$ is finite.
	Then the set of maximal elements is the minimal non-empty
	set w.r.t.~set inclusion that has non-empty
	intersection with each non-empty upper set.	
\end{proposition}
\begin{proof}
	Let~$T$ be a minimal non-empty upper set
	that has non-empty intersection with each non-empty upper set.
	
	Consider $x \in S$.
	By assumption, we have $y \in T \cap \upperset{\{x\}}$ for some $y \in X$.
	Since~$x$ is maximal, this gives $y \simeq x$.
	Then $x \in T$ since~$T$ is an upper set.	
	Hence $S \subseteq T$.
	
	Next, let $R$ be an upper set and $x \in R$. 
	By Lemma~\ref{prelim:general:order:lem:finite-maximal}, 
	there exists $y \in S$ such that $x \lesssim y$.
	Since~$R$ is an upper set, we have $y \in R \cap S$.
	Hence, by minimality of~$T$, we find $T \subseteq S$.
\end{proof}

\begin{corollary}
	Let~$X$ be finite and $S \subseteq X$ be a non-empty upper set.
	Then~$S$ contains an upper set that is an equivalence class.	
\end{corollary}
\begin{proof}
	By \cref{prelim:prop:maximal-elements-and-upper-sets} we 
	have $q \in S$ for some maximal element~$q \in X$.
	Then $[q]_{\simeq}$ is an upper set by \cref{prelim:general:order:cor:min-equiv-class}
	and we have $[q]_{\simeq} \subseteq S$.	
\end{proof} 

\subsection{Sequences of Upper and Lower Sets}
\label{prelim:general:order:sec:seq-upper-sets}

Here, we investigate sequences of increasing upper and lower sets
under certain finiteness conditions. We only give statements for upper sets,
since corresponding statements for lower sets follow analogously,
and we actually only need them for upper sets in later chapters.

\begin{proposition} \label{prelim:general:order:prop:seq-upper-sets}
	Let $S \subseteq Q$ be a proper upper set
	and $X \setminus S$ be finite.
	Then there exists a sequence of upper sets
	\[
	S_1 \subsetneq \cdots \subsetneq S_m = X
	\]
	with $S_1 = S$ and such that  $S_{i} \setminus S_{i-1}$ for $i \in \{2,\ldots,m\}$
	are all the equivalence classes in~$X \setminus S$.
\end{proposition}
\begin{proof}
	By finiteness, there exists a sequence
	of upper sets $S_1 \subsetneq \cdots \subsetneq S_m$
	with $S_1 = S$ and~$m$ maximal. Since $X$ is an upper set, we have $S_m = X$.
	
	First, assume there exists $i \in \{2,\ldots,m\}$ such that
	$S_{i} \setminus S_{i-1}$ is not contained in an equivalence
	class.
	By \cref{prelim:general:order:subset-in-equiv} there exists an upper set $T \subseteq X$
	with $( S_i \setminus S_{i-1} ) \cap T \ne \emptyset$
	and $( S_i \setminus S_{i-1} ) \not\subseteq T$.
	Thus, for the upper set $S_{i-1} \cup ( T  \cap S_i )$
	we have $S_{i-1} \subsetneq S_{i-1} \cup ( T  \cap S_i ) \subsetneq S_i$,
	which contradicts the maximal choice of~$m$.
	
	Next, let $x \in X \setminus S$. Then there exists a minimal $i \in \{2,\ldots,m\}$
	such that $x \in S_i$, so that $x \in S_i \setminus S_{i-1}$.
	Since $S_i \setminus S_{i-1}$ is contained in an equivalence
	class, we have $S_i \setminus S_{i-1} \subseteq [x]_{\simeq}$.
	By \cref{prelim:automata:cor:SsetminusT} 
	we find $S_i \setminus S_{i-1} = [x]_{\simeq}$.
\end{proof}

\begin{remark}
	\Cref{prelim:general:order:prop:seq-upper-sets} 
	can also be proved 
	by picking a proper upper set $S \subseteq S_{m-1}$ such that $|X \setminus S_{m-1}|$ has minimal size.
	Then~$S_{m-1}$ is a proper maximal upper set containing~$S$
	and by \cref{prelim:automata:prop:TsetminusS-strong-conn} 
	and  \cref{prelim:automata:cor:SsetminusT}
	the difference $X \setminus S_{m-1}$ equals an equivalence class.
	If $S \ne S_{m-1}$, inductively, 
	there exists a sequence of upper sets $S_1 \subsetneq \cdots \subseteq S_{m-1}$
	such that $S_i \setminus S_{i-1}$ for $i \in \{1,\ldots,m-1\}$
	are all the equivalence classes in $S_{m-1} \setminus S$.
	Now for each equivalence class $R \subseteq X$
	we either have $R \subseteq S_{m-1}$
	or $R \cap S_{m-1} = \emptyset$. Since the latter case implies $R \subseteq X \setminus S_{m-1}$,
	and then $R = X \setminus S_{m-1}$ since $X \setminus S_{m-1}$
	is an equivalence class, with $S_m = X$
	we also obtain a sequence as claimed.
\end{remark}

For finite~$X$, by picking for $S$ the empty set
or the set of maximal elements we obtain the following corollary
to \cref{prelim:general:order:prop:seq-upper-sets}.

\begin{corollary} \label{prelim:general:order:cor:seq-upper-sets}
	Suppose~$X$ is finite.
	\begin{enumerate} 
		\item  \label{prelim:general:order:cor:seq-upper-sets:1} There exists a sequence of upper sets
		\[
		\emptyset = S_0 \subsetneq S_1 \subsetneq \cdots \subsetneq S_m = X
		\]
		such that $S_{i} \setminus S_{i-1}$ for $i \in \{1,\ldots,m\}$
		are all the equivalence classes in~$X$.
		
		\item  \label{prelim:general:order:cor:seq-upper-sets:2} There exists a sequence of upper sets
		\[
		S_1 \subsetneq \cdots \subsetneq S_m = X
		\]
		such that $S_0$ equals the set of maximal elements
		and each equivalence class of a non-maximal element has the form
		$S_{i} \setminus S_{i-1}$ for some $i \in \{1,\ldots,m\}$.
	\end{enumerate} 
\end{corollary}

The sequences from \crefitem{prelim:general:order:cor:seq-upper-sets}{prelim:general:order:cor:seq-upper-sets:1}
are related to topological sortings in the following way.
Given a topological sorting $[x_1]_{\simeq}, \ldots, [x_m]_{\simeq}$ of the elements from $X / {\simeq}$ for finite~$X$,
we obtain a sequence as in \Crefitem{prelim:general:order:cor:seq-upper-sets}{prelim:general:order:cor:seq-upper-sets:1}
by setting $S_i = \upperset{\{x_1, \ldots, x_i\}}$ and $S_0 = \emptyset$.
Conversely, given a sequence as in \crefitem{prelim:general:order:cor:seq-upper-sets}{prelim:general:order:cor:seq-upper-sets:1},
the equivalence classes $S_i \setminus S_{i-1}$ for $i \in \{1,\ldots,m\}$
give a topological sorting of $X / {\simeq}$
since $x_i \lesssim x_j$ for $x_i \in S_i \setminus S_{i-1}$ and $x_j \in S_j \setminus S_{j-1}$ with $i \ne j$
implies $\upperset{\{x_i\}} \subsetneq \upperset{\{x_j\}}$ by \cref{prelim:eqn:reach-upperset}.
Thus, $\upperset{\{x_i\}} \subseteq S_i$
and $\upperset{\{x_j\}}\subseteq S_j$
give $S_i \subsetneq S_j$ and so $i \le j$.